% ---------------------------------------------------------------------------
% Results section (\input as 05-results from sigconf.tex).
%
% TODO: this file has eight result tables and their sourcing/methodology notes
% (as LaTeX comments, not yet body prose). The narrative text connecting them
% -- what to say about each table, in what order, with what emphasis -- still
% needs to be written; see the comment blocks above each table for the
% load-bearing claims that prose needs to state explicitly. In table order:
%   tab:pooling-ablation  -- does pooling (Table~\ref{tab:constructions}) beat
%                            individual training? Read this FIRST -- it
%                            justifies the pooled construction every other
%                            table in this section then assumes.
%   tab:results-2022/2023 -- the main model comparison (EXAMM/DLinear/Reversal/
%                            transformers), using the strategy defined right
%                            below. State explicitly: "do not read L/S 10 in
%                            isolation" ahead of tab:grid-robustness.
%   tab:sharpe             -- risk-adjusted view of the same cells.
%   tab:grid-robustness    -- the actual robustness finding, not L/S 10 alone.
%   tab:net-of-cost        -- same cells, after transaction costs.
%   tab:midhighmid         -- REPLICATION on a second universe: 4 random
%                            50-stock draws x 3 cohorts, 30-run ensembles.
%                            This is the only table that measures how much of a
%                            single-universe number is draw luck, and the answer
%                            is "a lot" (2023 IC spans -0.0195 to +0.0196 across
%                            draws). Prose must state that explicitly -- it
%                            qualifies every other table in this section.
%   tab:ina219_results     -- on-device (Raspberry Pi Zero) deployment cost,
%                            connects back to the parameter-count columns in
%                            tab:results-2022/2023.
%
% The \clearpage at the end of 04-datasets.tex already flushes anything
% queued from that section before this one's floats begin, so nothing
% drifts across the \input boundary. Do NOT add a second \clearpage here --
% two adjacent \clearpage calls force an extra, needless page break (this
% is what was leaving page 2 half-empty: the first one truncated section 4
% early, the second then stranded this section's heading alone on its own
% page before its tables could join it).
% ---------------------------------------------------------------------------

\section{Trading Results}
\label{sec:results}

% ---------------------------------------------------------------------------
% POOLING ABLATION -- goes FIRST in this section, ahead of the STRATEGY block
% below, because every other table in this section takes the pooled
% construction as given; this is the one table that actually justifies it
% rather than asserting it. L/S~10 here uses the identical conditional
% long/short strategy defined in the STRATEGY comment immediately below --
% not redefined here to avoid saying it twice.
%
% Does the pooled-panel construction (Table~\ref{tab:constructions}) actually
% outperform the individual (one-model-per-stock) alternative it replaces?
% Individual here is EXAMM trained one stock at a time (50 independent models x 10
% runs each, scripts/stock_run/anvil_cohort_individual.sb), same cohorts/windows as
% the pooled comparison below, predictions averaged across the 10 runs per stock
% (ensemble, matching how every other model in this paper is scored) before scoring
% IC/MSE/MAE and trading. Verified: expected_RET identical across all 10 runs per
% stock before averaging (rules out a file-pairing error); one run (WRB, cohort 2021)
% failed genome verification and was excluded, so that stock's ensemble is 9 runs.
%
% READING: MSE/MAE are nearly identical between individual and pooled in both years --
% pooling does not make each stock's point prediction more accurate. IC is where the
% gap actually lives (individual ~0 or negative both years; pooled clearly positive
% both years), which is the direct mechanism for the trading gap: a long/short
% strategy needs predictions that rank stocks against each other, and only the pooled
% model reliably produces that. State this explicitly in body prose -- it's the
% mechanistic explanation the rest of the section's numbers otherwise leave implicit.
% ---------------------------------------------------------------------------

Before comparing pooled EXAMM against other models, we first ask whether pooling itself
is doing any work. Table~\ref{tab:pooling-ablation} compares the pooled-panel construction
against training an independent EXAMM model per stock, holding the walk-forward cohorts,
10-run ensembling, and conditional long/short strategy fixed. Point-prediction accuracy
(MSE, MAE) is nearly identical between the two constructions in both years -- pooling does
not make any individual stock's forecast more accurate. The gap is entirely in
cross-sectional IC: individual training produces predictions with essentially no ranking
ability ($\text{IC}\approx 0$ in 2022, negative in 2023), while the pooled model is clearly
positive in both years. Because the long/short strategy trades on the relative ranking of
predicted returns across stocks rather than their absolute accuracy, this is the direct
mechanism behind pooling's trading advantage, and the reason every other table in this
section uses the pooled construction.

\begin{table*}[t]
\centering
\caption{Pooling vs. individual (one model per stock) training, both 10-run
ensembles, identical cohorts/windows/strategy as
Tables~\ref{tab:results-2022}--\ref{tab:results-2023}.}
\label{tab:pooling-ablation}
\begin{tabular}{llrrrr}
\toprule
Trade year & Construction & IC & MSE & MAE & L/S~10 \\
\midrule
2022 & Individual & $+0.0001$ & $4.726{\times}10^{-4}$ & $0.01584$ & $+13.36$ \\
2022 & \textbf{Pooled (ours)} & $\mathbf{+0.0030}$ & $\mathbf{4.630{\times}10^{-4}}$ & $\mathbf{0.01566}$ & $\mathbf{+14.40}$ \\
\midrule
2023 & Individual & $-0.0050$ & $3.037{\times}10^{-4}$ & $0.01232$ & $-1.08$ \\
2023 & \textbf{Pooled (ours)} & $\mathbf{+0.0195}$ & $\mathbf{2.930{\times}10^{-4}}$ & $\mathbf{0.01205}$ & $\mathbf{+31.99}$ \\
\bottomrule
\end{tabular}
\end{table*}

% ---------------------------------------------------------------------------
% STRATEGY: Lyu et al.'s own Algorithm 2 (arXiv 2410.17212), NOT an always-trade
% variant. Sort by predicted return; only trade at time t if ALL L candidate
% longs are predicted positive AND all S candidate shorts are predicted
% negative, else hold yesterday's book.
%
% ROW SET: FULL trading year, no truncation, for every model (250 days 2022,
% 249 days 2023) -- stock_pooled_loader.py now prepends the tail of
% validation as lookback context (leak-free -- val is fully in the past,
% used only as input, never as a target), so DLinear here already covers
% the identical full-year row set EXAMM does; no alignment trick needed.
%
% CONTROLS ARE LOAD-BEARING, NOT FILLER: Reversal (zero-parameter,
% predicted_RET = -RET[t]) is here because the source paper reports
% Algorithm 1 beating the market with no zero-training comparator. DLinear
% (42 fixed parameters, a 20-day trend/seasonal linear decomposition -- NOT
% evolved, not our design) is here because it is the smallest model tested
% that uses a genuinely different mechanism (multi-day trend-following vs.
% EXAMM's near-reversal policy -- daily rank-correlation with the Reversal
% control is 0.93-0.96 in every month of the test window). Without these
% rows this table would just restate the source paper's own unaudited claim.
% Reversal is a single deterministic fit; DLinear and EXAMM are true 10-seed
% ensembles.
%
% DO NOT READ L/S 10 ALONE AS THE FINDING. At a single book size, Reversal
% looks comparably "rescued" by the conditional filter, which would wrongly
% read as "the trading rule does all the work, not EXAMM." Swept across the
% full long-1..15 x short-1..15 grid (225 cells, same construction as the
% source paper's own Tables 5-6: results/trading/hybrid_grid_{2022,2023}
% {,_reversal,_dlinear}.csv), EXAMM is profitable AND beats both benchmarks
% in 93.8-100% of cells in BOTH years. Reversal and DLinear each collapse in
% at least one year (Reversal: 52% of cells positive in 2022; DLinear: only
% 17.8% beat the S&P in 2023 -- its trend-following edge does not survive
% into the choppier bull year). See Table~\ref{tab:grid-robustness}; state
% that finding in the body text, do not let L/S 10 imply it alone.
%
% Two parameter columns: every figure in a row (except Reversal) is a 10-run
% ensemble, so the per-model count alone understates the deployed model by
% 10x. DLinear's architecture is FIXED (not evolved), so its per-model count
% is exactly 42 every seed -- no range, unlike EXAMM.
%
% Benchmarks are rows, not columns: they are identical for all models within
% a year and have no book size. Net-of-cost is NOT computed yet for this
% (hybrid) strategy on this row set.
%
% SHARPE: computed for EVERY book size (5/10/15/20), not just L/S 10, from a
% full DAILY equity curve -- not reimplemented separately, to avoid the risk
% of a second implementation quietly diverging from the real one.
% scripts/stock_run/trade_portfolio_daily_curve.py pulls Portfolio.
% daily_hybrid_long_short_return's ACTUAL source via inspect.getsource,
% inserts exactly one line that appends CurrentCash + sum(share*price[t])
% (the same formula Stock.get_current_liquid uses) once per day, execs that
% patched source as a method on the real Portfolio class, then reruns the
% ORIGINAL unpatched method on a fresh Portfolio and asserts the two final
% returns match to 1e-6 before trusting the curve. All 24 cells here (EXAMM/
% DLinear/Reversal x 2022/2023 x 4 book sizes) validated to diff=0.000000pp.
% Sharpe is annualized (x sqrt(252)) on daily simple returns of the curve,
% rf=0. Given its own table (Table~\ref{tab:sharpe}) rather than packed into
% the L/S cells below -- two different-unit numbers per cell made both
% harder to scan.
%
% MARKET RETURNS -- SOURCES, VERIFIED, NOT ASSUMED
% --------------------------------------------------------------------------
% S&P 500 (-19.75% / +25.08%): NOT an external figure -- compounded from the
% "sprtrn" column already present in our own CRSP-sourced walk-forward data
% (trade_portfolio.py:benchmarks(), product of daily returns over rows
% 1..n_prices-2 of the trading window, exactly matching the window every
% other number in this table uses). Verified independently: fetched real
% S&P 500 (^GSPC) daily closes directly from Yahoo Finance's chart API
% (https://query1.finance.yahoo.com/v8/finance/chart/%5EGSPC), parsed the
% raw JSON with a script (NOT an AI-summarized read -- that path produced a
% wrong number for DJI below), and computed the simple close-to-close return
% for the same window (2022-01-03 -> 2022-12-29: 4796.56 -> 3849.28 =
% -19.7492%; 2023-01-03 -> 2023-12-28: 3824.14 -> 4783.35 = +25.0830%).
% Matches our compounded-sprtrn figures (-19.7496% / +25.0825%) to within
% 0.0004 percentage points.
%
% DJI (-9.20% / +13.80%): genuinely external -- no DJI column exists in our
% CRSP data. Source: Yahoo Finance chart API for ^DJI, raw JSON fetched with
% curl and parsed directly in Python -- explicitly NOT the AI-summarized
% read of the same JSON, which on a first attempt returned 32,981.55 for the
% 2022-12-30 close where the raw parse gives 33,147.25. That raw-parsed
% value was cross-checked against an independent source (StatMuse) and
% matches exactly. Close values used: 2022-01-03 = 36{,}585.06, 2022-12-29 =
% 33{,}220.80, 2023-01-03 = 33{,}136.37, 2023-12-28 = 37{,}710.10. Result is
% close to but not identical to the source paper's own reported DJI
% (-8.78% / +13.70%, their Table 4) -- same order of mismatch as the S&P
% discrepancy above, same reason: a different exact window boundary.
% ---------------------------------------------------------------------------

\begin{table*}[t]
\centering
\caption{Trading year 2022 (train $\le$2020, validate 2021), Algorithm
2~\cite{lyu2024neuroevolution}. Returns are annual \%, gross, 10-run
ensemble. \textbf{L/S~10} is pre-committed -- see
Table~\ref{tab:grid-robustness} before reading it alone, and
Table~\ref{tab:sharpe} for risk-adjusted returns.}
\label{tab:results-2022}
\resizebox{\textwidth}{!}{%
\begin{tabular}{lrrrrrrrrr}
\toprule
& \multicolumn{2}{c}{Parameters} & \multicolumn{3}{c}{Accuracy} & \multicolumn{4}{c}{Long/short book size} \\
\cmidrule(lr){2-3} \cmidrule(lr){4-6} \cmidrule(lr){7-10}
Model & Per model & Ensemble & IC & MSE & MAE & 5 & 10 & 15 & 20 \\
\midrule
\textbf{EXAMM} & $55$ & $549$ & $+0.0030$ & $4.630{\times}10^{-4}$ & $0.01566$
  & $+17.41$ & $\mathbf{+14.40}$ & $+17.83$ & $+16.39$ \\
DLinear & $42$ & $420$ & $+0.0116$ & $4.600{\times}10^{-4}$ & $0.01559$
  & $+9.72$ & $\mathbf{+5.39}$ & $-8.17$ & $-8.80$ \\
Reversal & $0$ & $0$ & $+0.0063$ & $9.280{\times}10^{-4}$ & $0.02244$
  & $-14.90$ & $\mathbf{+9.19}$ & $+5.71$ & $-2.41$ \\
Crossformer & $512{,}588$ & $5{,}125{,}880$ & $+0.0133$ & $4.537{\times}10^{-4}$ & $0.01548$
  & $+20.28$ & $\mathbf{+6.91}^{\dagger}$ & $0.00^{\ddagger}$ & $0.00^{\ddagger}$ \\
DeformTime  & $107{,}243$ & $1{,}072{,}430$ & --- & --- & --- & --- & --- & --- & --- \\
PatchTST    & $67{,}969$  & $679{,}690$     & --- & --- & --- & --- & --- & --- & --- \\
\midrule
\multicolumn{10}{l}{\footnotesize $^{\dagger}$gate opened on 2/250 days;
  $^{\ddagger}$gate opened on 0/250 days -- the book was never entered, so
  $0.00$ is ``never traded,'' not ``traded to a flat result.'' See
  Sec.~\ref{sec:gate-units} and Table~\ref{tab:gate-units}.} \\
\midrule
\multicolumn{6}{l}{Equal-weight buy-and-hold} & \multicolumn{4}{c}{$-13.99$} \\
\multicolumn{6}{l}{S\&P 500} & \multicolumn{4}{c}{$-19.75$} \\
\multicolumn{6}{l}{DJI} & \multicolumn{4}{c}{$-9.20$} \\
\bottomrule
\end{tabular}%
}
\end{table*}

\begin{table*}[b]
\centering
\caption{Trading year 2023 (train $\le$2021, validate 2022), Algorithm
2~\cite{lyu2024neuroevolution}. Returns are annual \%, gross, 10-run
ensemble. \textbf{L/S~10} is pre-committed -- see
Table~\ref{tab:grid-robustness} before reading it alone, and
Table~\ref{tab:sharpe} for risk-adjusted returns.}
\label{tab:results-2023}
\resizebox{\textwidth}{!}{%
\begin{tabular}{lrrrrrrrrr}
\toprule
& \multicolumn{2}{c}{Parameters} & \multicolumn{3}{c}{Accuracy} & \multicolumn{4}{c}{Long/short book size} \\
\cmidrule(lr){2-3} \cmidrule(lr){4-6} \cmidrule(lr){7-10}
Model & Per model & Ensemble & IC & MSE & MAE & 5 & 10 & 15 & 20 \\
\midrule
\textbf{EXAMM} & $77$ & $768$ & $+0.0195$ & $2.930{\times}10^{-4}$ & $0.01205$
  & $+41.77$ & $\mathbf{+31.99}$ & $+25.48$ & $+18.97$ \\
DLinear & $42$ & $420$ & $+0.0123$ & $2.940{\times}10^{-4}$ & $0.01209$
  & $+19.17$ & $\mathbf{+7.84}$ & $+1.78$ & $-4.40$ \\
Reversal & $0$ & $0$ & $+0.0249$ & $5.820{\times}10^{-4}$ & $0.01734$
  & $+63.10$ & $\mathbf{+14.83}$ & $+15.07$ & $-1.45$ \\
Crossformer & $512{,}588$ & $5{,}125{,}880$ & $+0.0213$ & $2.924{\times}10^{-4}$ & $0.01203$
  & $+38.12$ & $\mathbf{+16.40}$ & $+8.58$ & $-1.04$ \\
DeformTime  & $107{,}243$ & $1{,}072{,}430$ & --- & --- & --- & --- & --- & --- & --- \\
PatchTST    & $67{,}969$  & $679{,}690$     & --- & --- & --- & --- & --- & --- & --- \\
\midrule
\multicolumn{6}{l}{Equal-weight buy-and-hold} & \multicolumn{4}{c}{$+7.91$} \\
\multicolumn{6}{l}{S\&P 500} & \multicolumn{4}{c}{$+25.08$} \\
\multicolumn{6}{l}{DJI} & \multicolumn{4}{c}{$+13.80$} \\
\bottomrule
\end{tabular}%
}
\end{table*}

% No \clearpage here. The wide pair above are pinned [t]/[b], so LaTeX's
% two-column algorithm already places table* floats before same-column
% ([t]/[b]) tables that follow them in the source -- a table* cannot be
% pushed earlier than its source position, only later, so it cannot land
% mid-text. Forcing a page break here was what stranded the wide pair alone
% on one page and pushed Sharpe/grid-robustness onto their own near-empty
% page 3; removing it lets the single-column tables fill whatever column
% space the wide pair's page leaves over instead of abandoning it.

\begin{table*}[t]
\centering
\caption{Annualized Sharpe ratio ($r_f{=}0$) by book size, same data as
Tables~\ref{tab:results-2022}--\ref{tab:results-2023}. \textbf{L/S~10}
bolded to match; not yet computed for Crossformer/DeformTime/PatchTST.}
\label{tab:sharpe}
\begin{tabular}{llrrrr}
\toprule
Trade year & Model & 5 & 10 & 15 & 20 \\
\midrule
2022 & \textbf{EXAMM} & $0.90$ & $\mathbf{0.98}$ & $1.36$ & $1.49$ \\
2022 & DLinear         & $0.55$ & $\mathbf{0.41}$ & $-0.51$ & $-0.64$ \\
2022 & Reversal        & $-0.62$ & $\mathbf{0.56}$ & $0.43$ & $-0.12$ \\
\midrule
2023 & \textbf{EXAMM} & $2.04$ & $\mathbf{2.32}$ & $2.39$ & $2.07$ \\
2023 & DLinear         & $0.99$ & $\mathbf{0.64}$ & $0.22$ & $-0.41$ \\
2023 & Reversal        & $2.53$ & $\mathbf{1.04}$ & $1.25$ & $-0.10$ \\
\bottomrule
\end{tabular}
\end{table*}

% ---------------------------------------------------------------------------
% Grid-robustness note -- state this (or a summary of it) in body text next
% to the tables above. It is the actual finding; the L/S 10 column alone is
% not, since Reversal looks comparably "rescued" at that one cell and DLinear
% looks fine at that one cell despite its 2023 grid collapsing.
% Source: results/trading/hybrid_grid_{2022,2023}{,_reversal,_dlinear}.csv,
% each a 225-cell sweep (long 1..15 x short 1..15), matching the construction
% of Lyu et al.'s own Tables 5-6.
% ---------------------------------------------------------------------------
\begin{table*}[b]
\centering
\caption{Robustness across the full long-1..15 $\times$ short-1..15 book-size
grid (225 cells), Algorithm 2. Percent of cells that are profitable / beat
equal-weight buy-and-hold / beat the S\&P 500. EXAMM is the only model that
clears $>$90\% on every measure in both years.}
\label{tab:grid-robustness}
\begin{tabular}{llrrr}
\toprule
Trade year & Model & Cells $>0$ & Beat B\&H & Beat S\&P \\
\midrule
2022 & \textbf{EXAMM} & $\mathbf{93.8\%}$ & $\mathbf{99.6\%}$ & $\mathbf{99.6\%}$ \\
2022 & DLinear         & $74.2\%$ & $98.7\%$ & $99.6\%$ \\
2022 & Reversal        & $52.0\%$ & $83.1\%$ & $89.3\%$ \\
\midrule
2023 & \textbf{EXAMM} & $\mathbf{100.0\%}$ & $\mathbf{100.0\%}$ & $93.3\%$ \\
2023 & DLinear         & $71.1\%$ & $52.9\%$ & $17.8\%$ \\
2023 & Reversal        & $96.4\%$ & $93.3\%$ & $51.1\%$ \\
\bottomrule
\end{tabular}
\end{table*}

% ---------------------------------------------------------------------------
% GATE UNITS. This subsection exists because Crossformer's two 0.00 cells in
% Table 2 are NOT a forecasting result and must not be read as one.
%
% Algorithm 2 opens the book only if all L longs are predicted positive and
% all S shorts negative. That zero is an ABSOLUTE constant compared against
% raw model output, so the opening rate is governed by the ratio
% (prediction mean)/(prediction sd) -- i.e. by how far each model shrank
% toward the unconditional mean under MSE -- and not by forecast skill.
% Measured on the 2022 test cross-sections (10-seed ensembles):
%     EXAMM        83.4% of predictions > 0, pooled sd 2.13e-03
%     Crossformer  99.3% of predictions > 0, pooled sd 3.83e-04  (5.6x shrunk)
% Crossformer's cloud therefore sits ~2.8 sd above zero and its short leg
% essentially cannot clear it: 2/250 openings at L/S 10, 0/250 at L/S 15
% and 20. Same model, 2023: 116/249 at L/S 10 -- healthy. Nothing about the
% architecture changed between those two rows; only the output scale did.
%
% WHY NOT JUST RECALIBRATE THE PREDICTIONS. The textbook fix is a
% validation-fit affine map pred' = a + b*pred restoring validation-period
% return dispersion. Measured, it is NOT usable here: EXAMM's own b is 16.63
% and applying it moves EXAMM 2022 L/S 10 from +14.40 to +0.39, and makes the
% book-size profile wildly unstable (+37.11/+0.39/-9.36/+0.74 at L/S
% 5/10/15/20 vs a flat +17.41/+14.40/+17.83/+16.39 uncalibrated). A transform
% that rewrites every model's headline number is not a fix, it is a different
% experiment.
%
% WHAT WE DO INSTEAD. Move the threshold, not the predictions: subtract a
% per-model constant tau_q = the q-th percentile of that model's own pooled
% predictions. A CONSTANT shift is rank-preserving within every day, so the
% long/short SELECTION -- and therefore IC, MSE, MAE and Sharpe -- is
% bit-identical; only the gate moves. Published Algorithm 2 is exactly the
% tau=0 point of each model's own curve, so no number in Tables 2-3 changes.
% The diagnosis is that tau=0 lands at completely different places per model:
% q = 15.1% for EXAMM but 0.5% for Crossformer in 2022 (56.5% / 23.1% in
% 2023). Matching q is what makes the row comparable.
%
% Source: scratchpad gate_threshold_sweep.py + trade_portfolio.py
% (--strategy daily_hybrid_long_short_return), same windows and same 10-seed
% ensembles as Tables 2-3.
% ---------------------------------------------------------------------------
\subsection{The gate threshold is not in comparable units}
\label{sec:gate-units}

Algorithm 2 opens the book only when all $L$ longs are predicted positive and
all $S$ shorts negative. That zero is an \emph{absolute} constant compared
against raw model output, so how often the gate opens depends on how far a
model shrank toward the unconditional mean under MSE -- not on how well it
forecasts. In 2022 Crossformer put $99.3\%$ of its predictions above zero with
a pooled standard deviation $5.6\times$ smaller than EXAMM's, placing its
entire cross-section roughly $2.8$ standard deviations above the threshold.
Its short leg then cannot clear zero: the gate opened on $2/250$ days at
L/S~10 and on $0/250$ days at L/S~15 and~20. The same model in 2023 opened on
$116/249$ days. The architecture did not change between those rows; only the
output scale did.

We therefore move the threshold rather than the predictions, subtracting from
each model a constant $\tau_q$ equal to the $q$-th percentile of its own
pooled predictions. A constant shift preserves within-day ranks, so the
long/short selection -- and with it IC, MSE, MAE and Sharpe -- is unchanged;
only the gate moves. Published Algorithm 2 is the $\tau{=}0$ point of each
model's own curve, so every number in
Tables~\ref{tab:results-2022}--\ref{tab:results-2023} stands as reported. The
defect is that $\tau{=}0$ sits at $q{=}15.1\%$ for EXAMM but $q{=}0.5\%$ for
Crossformer, so the two were never gated at comparable points of their own
distributions.

Table~\ref{tab:gate-units} sweeps $q$ jointly. Crossformer's two degenerate
cells become well-defined once it is allowed to trade, and in 2022 EXAMM leads
it at every threshold in the profitable region. We report this as a
\emph{diagnostic and not as a corrected result}, for two reasons. There is no
principled value of $q$: the only non-arbitrary threshold is zero applied to
calibrated predictions, which is the affine map rejected above, and every other
choice is a free parameter. And $\tau_q$ is a percentile of the pooled
\emph{test} predictions -- no labels enter, but test-period information does,
which is admissible when the whole curve is shown and no point is selected, and
not admissible for a headline number.
Tables~\ref{tab:results-2022}--\ref{tab:results-2023} therefore remain the
reported result, at $\tau{=}0$,
exactly as Algorithm 2 defines it. A single calibrated column would require
$\tau$ fit on validation predictions and carried forward, which we could not
reconstruct for the transformer baselines.

Two further things should be read honestly alongside the sweep. First,
Crossformer is the more \emph{accurate} model
in both years ($+0.0133$ vs.\ $+0.0030$ IC in 2022; $+0.0213$ vs.\ $+0.0195$
in 2023) -- consistent with DLinear also out-predicting EXAMM in 2022, and
with our finding that Algorithm 2's return is driven more by when the gate
opens than by cross-sectional ranking quality. Second, the 2023 sweep has no
stable structure for any model, and EXAMM's own $+31.99$ sits at $q{=}56.5\%$
where the neighbouring grid point returns $+14.67$. The gate is a market-timing
filter, and in a year it times badly no model's number is threshold-robust.

\begin{table*}[t]
\centering
\caption{Algorithm 2 with the gate threshold set at the $q$-th percentile of
each model's \emph{own} pooled predictions (annual \%, gross, 10-run
ensemble, L/S~10). Rank-preserving, so IC/MSE/MAE/Sharpe are unchanged from
Tables~\ref{tab:results-2022}--\ref{tab:results-2023}. Crossformer's 2022
cells that were $0.00$ at $\tau{=}0$ are shown at L/S~15 for contrast.}
\label{tab:gate-units}
\begin{tabular}{llrrrrrrrr}
\toprule
Trade year & Model & $q{=}2$ & $5$ & $10$ & $15$ & $20$ & $25$ & $30$ & $50$ \\
\midrule
2022 & \textbf{EXAMM} & $+35.54$ & $+30.02$ & $+25.79$ & $+14.43$ & $+15.52$ & $+0.01$ & $-10.82$ & $-10.29$ \\
2022 & DLinear         & $+0.31$ & $-7.54$ & $-0.67$ & $+18.96$ & $+6.59$ & $+0.52$ & $+9.46$ & $+5.62$ \\
2022 & Crossformer     & $+21.79$ & $+18.16$ & $+7.72$ & $-23.17$ & $-1.31$ & $+1.04$ & $-3.65$ & $-8.36$ \\
2022 & Reversal        & $-11.07$ & $+41.90$ & $+39.58$ & $+0.21$ & $+2.78$ & $-8.79$ & $+6.09$ & $+7.05$ \\
\midrule
2023 & \textbf{EXAMM} & $-3.39$ & $+7.58$ & $+11.86$ & $+11.02$ & $-2.59$ & $+8.47$ & $+20.47$ & $+14.67$ \\
2023 & DLinear         & $-19.48$ & $-13.06$ & $-34.95$ & $-22.45$ & $-11.00$ & $-18.79$ & $-13.97$ & $+2.64$ \\
2023 & Crossformer     & $-0.69$ & $+3.41$ & $+0.99$ & $+6.39$ & $-14.22$ & $+12.16$ & $+21.92$ & $+24.93$ \\
2023 & Reversal        & $-17.91$ & $-18.74$ & $-21.13$ & $+8.51$ & $+21.74$ & $+16.10$ & $+15.67$ & $+15.05$ \\
\midrule
\multicolumn{10}{l}{\emph{2022 at L/S~15, where Crossformer returned $0.00$ under
  $\tau{=}0$:}} \\
2022 & \textbf{EXAMM} & $+4.63$ & $+37.72$ & $+7.55$ & $+18.01$ & $+2.38$ & --- & --- & --- \\
2022 & Crossformer     & $-1.86$ & $+17.45$ & $+11.12$ & $+17.02$ & $-11.15$ & --- & --- & --- \\
\bottomrule
\end{tabular}
\end{table*}

% ---------------------------------------------------------------------------
% NET OF TRANSACTION COSTS -- source, and a real bug found+fixed to get here.
% --------------------------------------------------------------------------
% --use-tc charges the per-share CRSP TRAN_COST on every buy/sell/short/cover, via
% Financial_toolbox's Stock class (an EXTERNAL dependency, ~/Financial_toolbox,
% zimenglyu/Financial_toolbox -- not part of this repo).
%
% BUG FOUND (2026-07-31): on the short side, portfolio.py credited the flat NOMINAL
% quota (self.current_cash_amount += quota_per_stock) regardless of TC, while
% Stock.short_stock already reduces share count for the fee (shares = quota/(price+TC)).
% The open leg thus ignored the fee entirely while the close leg paid it twice (fewer
% shares, then TC again on cover) -- net effect: enabling TC could INCREASE a strategy's
% return, which is not a real possibility for an actual friction cost. Confirmed both
% algebraically (a flat-price short round trip showed a phantom gain instead of a
% breakeven-minus-fee loss) and empirically: before the fix, ALL 8 DLinear cells and
% 2023 Reversal improved under --use-tc; EXAMM's 8 cells happened to stay directionally
% correct (net <= gross) but that looks incidental to this book's trade mix, not a
% property the old code guaranteed. Bug was present identically in all 4 of the
% toolbox's short-selling strategies (portfolio_shorting_return, daily_long_short_return,
% daily_hybrid_long_short_return, long_short_return), not just the one used here.
%
% FIX: Stock.short_stock now returns the actual realized credit -- shares_shorted *
% stock_price[time] (post-TC share count, but valued at the RAW price, matching what the
% long side already did implicitly) -- instead of callers assuming the nominal quota was
% received in full. Equals the nominal quota exactly (bit-identical) when TC is off or
% trading at bid/ask, so the no-TC path is untouched. VERIFIED: all 6 gross L/S-10 numbers
% in Tables~\ref{tab:results-2022}--\ref{tab:results-2023} reproduced exactly post-fix
% (+14.40/+5.39/+9.19/+31.99/+7.84/+14.83), and all 24 cells below now show net <= gross.
%
% NOT YET PUSHED to zimenglyu/Financial_toolbox (origin/main) -- this fix currently only
% exists in the local ~/Financial_toolbox checkout. Tell Dr. Lyu before citing these
% numbers as reproducible from a fresh clone of her repo.
% ---------------------------------------------------------------------------
\begin{table*}[t]
\centering
\caption{Net of transaction costs (per-share CRSP \texttt{TRAN\_COST}, charged on every
trade), same construction and book sizes as
Tables~\ref{tab:results-2022}--\ref{tab:results-2023}. \textbf{L/S~10} bolded to match;
not yet computed for Crossformer/DeformTime/PatchTST.}
\label{tab:net-of-cost}
\begin{tabular}{llrrrr}
\toprule
Trade year & Model & 5 & 10 & 15 & 20 \\
\midrule
2022 & \textbf{EXAMM} & $+14.21$ & $\mathbf{+12.73}$ & $+16.72$ & $+15.82$ \\
2022 & DLinear         & $+5.59$  & $\mathbf{+3.13}$  & $-9.36$  & $-9.28$  \\
2022 & Reversal        & $-19.23$ & $\mathbf{+5.47}$  & $+3.49$  & $-3.40$  \\
\midrule
2023 & \textbf{EXAMM} & $+39.14$ & $\mathbf{+30.24}$ & $+24.34$ & $+18.36$ \\
2023 & DLinear         & $+15.19$ & $\mathbf{+5.25}$  & $+0.13$  & $-5.11$  \\
2023 & Reversal        & $+57.55$ & $\mathbf{+11.68}$ & $+12.80$ & $-2.38$  \\
\bottomrule
\end{tabular}
\end{table*}

% ---------------------------------------------------------------------------
% MID/HIGH-MID REPLICATION -- 4 random 50-stock draws x 3 walk-forward cohorts.
%
% WHY THIS TABLE EXISTS: every other number in this section comes from ONE
% 50-stock universe. Four independent draws from a 200-stock pool make the
% sampling variability measurable for the first time, and it is large: in trade
% year 2023 the draws span IC -0.0195 to +0.0196 -- wider than the effect being
% measured. A single-universe result cannot separate method from draw.
%
% ENSEMBLE: 30 runs = seeds 1-30 (three 10-run batches pooled by averaging
% predictions, which for equal batch sizes is exactly a 30-run mean).
% Campaign dirs: <set>_<cohort>{,_rep2,_rep3} -> combined as *_e123.
%
% ROBUSTNESS OF THE ENSEMBLE SIZE (measured, 480 total training runs; seeds
% 31-40 were also run): mean IC is FLAT in ensemble size -- 10-seed +0.0056,
% 20-seed +0.0056, 30-seed +0.0057, 40-seed +0.0058 -- so the accuracy numbers
% here do not depend on this choice. Trading is NOT monotone in ensemble size
% (30-seed +16.63% vs 40-seed +12.97%) and the six exchangeable 20-seed pairings
% span +6.90% to +17.81%, i.e. ~11 pp of pure combination noise. Treat the
% trading column as carrying roughly +/-5 pp of seed noise; the IC column is the
% stable one (per-run sd ~0.0047 from 72 paired comparisons).
%
% BENCHMARK CHOICE: EW buy-and-hold is primary because it holds universe, period,
% weighting and survivorship fixed. The S&P 500 does NOT share the universe --
% these mid/high-mid caps trailed the cap-weighted index by 9.7 / 13.9 / 17.7 pp
% in 2022/23/24 -- so an S&P comparison confounds strategy with the universe's
% size tilt, a confound larger than the effect measured. Kept as context only.
% (Lyu et al. report B&H, DJI and S&P; retaining B&H + S&P matches the source,
% but the claim is made against B&H.)
%
% Scorecard: beats EW B&H 9/12, beats S&P 4/12, positive 11/12. All four 2022
% cells beat B&H; the two bull years go 5/8. Note 2022's mean (+27.54%) is
% carried by set2 (+58.35%) and set3 (+38.99%) -- set4 is NEGATIVE (-7.64%) that
% year despite the highest IC in the column (+0.0251), which is the clearest
% single illustration that IC and one-book-size trading return decouple.
%
% Sources: test_output/mse_pooled_mid_highmid/<set>_<cohort>_e123/ensemble_test/,
% data datasets/walkforward/mid_highmid_price/<set>/<cohort>_aligned/.
% Build: scripts/stock_run/{clean_mid_highmid.py, build_mid_highmid_splits.sh,
% anvil_mse_pooled_mid_highmid.sb (SEED_OFFSET/RUN_TAG), eval_mid_highmid.sh}.
% ---------------------------------------------------------------------------
\begin{table*}[t]
\centering
\caption{Replication on four independent 50-stock draws from a 200-stock mid/high-mid-cap CRSP universe (2\,$\le$\,$M$\,$<$\,50\,USD\,bn, 20+ years continuous history), \textbf{30-run ensemble}. Trading is Algorithm~2~\cite{lyu2024neuroevolution} at L/S~10, annual \%. Equal-weight buy-and-hold shares the universe, period, and survivorship filter and is the primary comparator; the S\&P~500 does not and is reported for context.}
\label{tab:midhighmid}
\resizebox{\textwidth}{!}{%
\begin{tabular}{llrrrrrrr}
\toprule
& & \multicolumn{3}{c}{Accuracy} & \multicolumn{2}{c}{Trading (\%)} & \multicolumn{2}{c}{Benchmarks (\%)} \\
\cmidrule(lr){3-5} \cmidrule(lr){6-7} \cmidrule(lr){8-9}
Draw & Trade yr & IC & MSE & MAE & Gross & Net of TC & EW B\&H & S\&P 500 \\
\midrule
set1 & 2022 & $-0.0013$ & $5.425{\times}10^{-4}$ & $0.01674$ & $+20.48$ & $+19.87$ & $-7.14$ & $-19.75$ \\
set1 & 2023 & $+0.0136$ & $3.218{\times}10^{-4}$ & $0.01240$ & $+10.89$ & $+9.71$ & $+5.21$ & $+25.08$ \\
set1 & 2024 & $+0.0055$ & $3.566{\times}10^{-4}$ & $0.01234$ & $+11.03$ & $+9.81$ & $-4.70$ & $+24.54$ \\
\addlinespace
set2 & 2022 & $+0.0136$ & $5.583{\times}10^{-4}$ & $0.01737$ & $+58.35$ & $+57.02$ & $-11.39$ & $-19.75$ \\
set2 & 2023 & $+0.0196$ & $3.456{\times}10^{-4}$ & $0.01298$ & $+18.93$ & $+17.05$ & $+15.41$ & $+25.08$ \\
set2 & 2024 & $+0.0119$ & $3.052{\times}10^{-4}$ & $0.01197$ & $+13.35$ & $+12.99$ & $+12.17$ & $+24.54$ \\
\addlinespace
set3 & 2022 & $+0.0086$ & $5.668{\times}10^{-4}$ & $0.01707$ & $+38.99$ & $+38.41$ & $-11.36$ & $-19.75$ \\
set3 & 2023 & $-0.0195$ & $3.295{\times}10^{-4}$ & $0.01281$ & $+10.52$ & $+9.75$ & $+14.80$ & $+25.08$ \\
set3 & 2024 & $-0.0117$ & $3.292{\times}10^{-4}$ & $0.01229$ & $+9.64$ & $+7.68$ & $+13.15$ & $+24.54$ \\
\addlinespace
set4 & 2022 & $+0.0251$ & $4.361{\times}10^{-4}$ & $0.01533$ & $-7.64$ & $-7.87$ & $-10.22$ & $-19.75$ \\
set4 & 2023 & $-0.0005$ & $2.424{\times}10^{-4}$ & $0.01116$ & $+2.42$ & $+2.03$ & $+9.45$ & $+25.08$ \\
set4 & 2024 & $+0.0033$ & $2.337{\times}10^{-4}$ & $0.01047$ & $+12.56$ & $+12.46$ & $+6.72$ & $+24.54$ \\
\midrule
\multicolumn{2}{l}{\textbf{Mean 2022}} & $\mathbf{+0.0115}$ & & & $\mathbf{+27.54}$ & $+26.86$ & $-10.03$ & $-19.75$ \\
\multicolumn{2}{l}{\textbf{Mean 2023}} & $\mathbf{+0.0033}$ & & & $\mathbf{+10.69}$ & $+9.64$ & $+11.22$ & $+25.08$ \\
\multicolumn{2}{l}{\textbf{Mean 2024}} & $\mathbf{+0.0022}$ & & & $\mathbf{+11.64}$ & $+10.73$ & $+6.83$ & $+24.54$ \\
\multicolumn{2}{l}{\textbf{Mean, all 12}} & $\mathbf{+0.0057}$ & & & $\mathbf{+16.63}$ & $+15.74$ & $+2.67$ & $+9.96$ \\
\bottomrule
\end{tabular}%
}
\end{table*}

% No trailing \clearpage: there is no \input after this one in sigconf.tex
% yet, so a forced break here only wastes a page. If a new section is added
% after this one, put ONE \clearpage at the START of that new file (matching
% the single boundary-flush pattern used at the top of this file), not here.



\subsection{Raspberry Pi Zero Results}
\label{sec:ina-219}

\begin{table*}[t]
\centering
\footnotesize
\begin{tabular}{lrr}
\toprule
Metric & 2020 & 2021 \\
\midrule
Params              & 55      & 77     \\
Voltage (V)         & 4.72    & 4.71   \\
Current (mA)        & 337.1   & 345.2  \\
Power (mW)          & 1586.6  & 1622.4 \\
Idle Pwr (mW)       & 1131.1  & 1124.5 \\
Gross E/file (mJ)   & 7.61    & 4.86   \\
Net E/file (mJ)     & 2.18    & 1.48   \\
Latency ($\mu$s)    & 9.59    & 12.04  \\
Throughput (k/s)    & 113.3   & 87.7   \\
\bottomrule
\end{tabular}
\caption{Model size, INA219 power/energy, and latency/throughput, averaged across 10 runs per year on Raspberry Pi Zero.}
\label{tab:ina219_results}
\end{table*}   